% borrow-rate.tex — Utilization-based borrow rate (JumpRate) + separate funding
% Kontrol: prove_ins_borrowRate*, prove_ins_utilizationBound*

\section{Borrow Rate}\label{sec:borrow-rate}

The PLP (borrower) pays two separate rates to the underwriter (lender):
a unidirectional borrow rate (always PLP $\to$ underwriter) and a bidirectional
funding rate (\cref{sec:funding-rate}). These are economically distinct mechanisms
with separate accumulators, following the Primer on Perpetuals
\cite{angeris2022perpetuals}.

\subsection{Borrow Rate: JumpRate Curve}

\begin{definition}[Borrow Rate]\label{def:borrow-rate}
\begin{equation}\label{eq:borrow-rate}
  \rborrow(\Ua) = \premfactor + \mlow \cdot \min(\Ua, \kink)
                 + \mhigh \cdot \max(0,\, \Ua - \kink)
\end{equation}
where:
\begin{itemize}
  \item $\premfactor > 0$: base premium (floor), calibrated from econometric WTP;
  \item $\Ua = L_{\text{borrowed}} / L_{\text{total}} \in [0, 1]$: utilization ratio;
  \item $\kink \in (0, 1)$: target utilization (e.g., $\kink = 0.8$);
  \item $\mlow > 0$: slope below kink (gentle increase);
  \item $\mhigh \gg \mlow$: slope above kink (jump multiplier, strong deterrent).
\end{itemize}
\end{definition}

\begin{proposition}[Borrow Rate Properties]\label{prop:borrow-properties}
\begin{enumerate}[label=(\roman*)]
  \item $\rborrow \geq \premfactor > 0$ for all $\Ua \geq 0$ (underwriters always earn at least $\premfactor$).
  \item $\rborrow$ is continuous and piecewise-linear in $\Ua$.
  \item $\rborrow$ is non-decreasing in $\Ua$ (higher utilization $\to$ higher borrow cost).
  \item At $\Ua = 1$: $\rborrow = \premfactor + \mlow \kink + \mhigh(1 - \kink)$
    (very expensive, incentivizing new underwriter deposits or PLP exit).
\end{enumerate}
\end{proposition}

\subsection{Borrow Rate Accrual}

\begin{definition}[Borrow Fee Accumulator]\label{def:borrow-accumulator}
At each hook callback (afterSwap, afterAddLiquidity, afterRemoveLiquidity),
the borrow fee growth accumulator updates:
\begin{equation}\label{eq:borrow-accumulator}
  g^{\text{borrow}} \mathrel{+}= \rborrow(\Ua) \cdot \Delta t
\end{equation}
where $\Delta t$ is the time elapsed since the last callback (in blocks or seconds).
This accumulator is global (per insurance pool) and distributed to underwriters
pro-rata by liquidity share, using the Synthetix staking pattern
\cite{batog2018scalable}.
\end{definition}

\subsection{Separate Accounting Principle}

\begin{remark}[Two Accumulators, Two Purposes]
The borrow rate and funding rate (\cref{sec:funding-rate}) use separate
accumulators:
\begin{itemize}
  \item $g^{\text{borrow}}$: always positive, PLP $\to$ underwriter.
    Compensates capital lockup and tail risk.
  \item $g^{\text{funding}}$: bidirectional. Corrects mark-index mispricing.
    When positive, PLP pays; when negative, underwriter pays.
\end{itemize}
The Primer on Perpetuals \cite[Section~3]{angeris2022perpetuals} establishes that
conflating these mechanisms breaks funding convergence: a funding term that
cannot go negative (because it's floored by the borrow rate) cannot fully
correct underpricing.
\end{remark}
